\subsection{Critical Technical Practice}
Whwn author Philip Agre calls for a "critical technical practice" it is not a simple proposal to add critique \emph{after} the fact. What is meant, is to effectively braid critique and critical thought into the “daily work of trying to get things built and working” so that technical design is continually framed and reframed by reflective inquiry \parencite[p.~155]{Agre_1998}. Crucially, Agre argues that the designer must cultivate a “split identity—one foot planted in the craft work of design and the other foot planted in the reflexive work of critique,” using historical and sociotechnical understanding to choose problems, diagnose impasses, and motivate alternatives from within practice \parencite[p.~155]{Agre_1998}. Read closely, this is a methodological stance about where knowledge lives: not in external commentary upon technology, but in an iterative, situated process that surfaces how “vague or ambiguous language,” notational habits, and default assumptions can bury conceptual trouble and reproduce a field’s narrow self-understanding \parencite[p.~155]{Agre_1998}. For audio toolmaking, Agre’s lesson is specific: defaults, interface metaphors, and algorithmic parameters are hypotheses about the human user and sonic world; a critical technical practice makes those hypotheses visible, testable, and revisable in the midst of building.

Phoebe Sengers’ program of reflective design extends Agre’s stance by putting to work \emph{reflexivity} as a design value. One that can transform technologies in addition to evaluating them critically. If we were to briefly compare the two approaches presented here: Agre emphasizes the structural need to hold together craft and critique, while Sengers et al., similarly, shows how to embed reflection into methods (provotyping, value scenarios, design fictions, and critical documentation). Within the "Principles of Reflective Design," these methods state the affordances and assumptions of artifacts are elicited, contested, and re-specified through the act of making experience \parencite[~55]{Sengers_2005}. The conversation between Agre and Sengers thus moves, in a way, from posture to practice: Agre demands that we refuse the closure of technical framings; Sengers provides techniques that open framings in concrete, situated ways. From a more audio-focused context, reflective design means interrogating the “naturalness” of, say, equal-tempered scales, studio-grounded genre presets, or linear time metaphors in DAW timelines. The subsequent reflective design practice would mean building alternative logics as fundamental technical proposals, not as an form of afterthought to be applied after the fact.

Looking at both of these approaches together, critical technical practice provides the methodological backbone for Sonic Speculocultural Technopoiesis. We will hold to Agre’s border-spanning stance as a discipline of practice—pairing each implementation task with a reflexive moment or phase aimed at surfacing hidden commitments. In direct addition we adopt Sengers’ repertoire to actualize those moments in the workflow: issue templates that query values, code comments that narrate design rationales, critique check-ins with practitioners, and probes that solicit counter-examples from Black sonic traditions \parencite{Agre_1998,Sengers_2005}. On Agre’s own terms, anticipated “impasses” (for instance, when a model cannot account for polyrhythmic feel or timbral grain as social practice) will be treated as productive signals to reframe problems and generate alternatives from within the build \parencite[p.154-155]{Agre_1998}. This is the mode by which critique feeds back into technopoiesis: reflection becomes material in code, and in interface, and in signal flow.

\subsection{Research-Creation and Practice-Based Knowledge}
In teh foundational text, "How to Make Art at the End of the World," author Natalie Loveless articulates research-creation as a university-sited practice in which “the artistic” is mobilized as a mode of inquiry attentive to how “form makes worlds,” working micropolitically “bit by bit” to reshape knowledge-making from the inside \parencite[p.~102]{Loveless_2019}. For Loveless, as purely in concer wit the fundamentals of research-creation, theory is to \emph{applied} practice, theory is \emph{generated} through practice(making). Research-creation is driven by "cathexis" and situated curiosity that often outstrips what existing disciplinary frameworks can justify \parencite[p.~15]{Loveless_2019}. In a closer reading we can extract two methodological provocations germane to sonic tool design: firstly, that the artifact is not a demonstration of prior knowledge but the site in which concepts are invented; and secondly, that institutional forms and systems must need to accommodate this invention, because the criteria of rigor resides in the coherence and care of the world-making that is performed \parencite[p.54-56]{Loveless_2019}. For Sonic Speculocultural Technopoiesis, this means treating plugins, interfaces, and datasets as argumentative forms that enact claims about time, rhythm, listening, identity, and community. The validity of the tool is inseparable from the ways they compose (\emph{human-world}) relations in both practice and performance.

The authors Erin Manning and Brian Massumi rendezvous with Loveless while pushing the ontology of practice further: research-creation, they argue, operates at a “prebifurcation” level where making is already "thinking-in-action" and conceptualization is a practice in its own right \parencite[p.~89]{Manning_Massumi_2014}. Technique, on their account, is processual and immanent; the task is to “construct the conditions for a speculative pragmatism” in which results are not preprogrammed but emerge as events catalyzed by initial constraints and attunements \parencite[p.~89]{Manning_Massumi_2014}. This shifts methodological emphasis from after-the-fact evaluation to before-the-fact conditioning: setting up collaborative situations and material affordances such that new conceptual articulations can occur. In audio design, this could potentially entail building instruments that are less interface-to-function mappings and operate more like “event conditioners.” These would serve as systems tuned to elicit feel and response that cannot be specified in advance, but more like iteratively stabilized through use.

Bringing Loveless into conversation with Manning and Massumi yields a concrete research protocol. Research-creation supplies the institutional and ethical frame (form as world-making, care as rigor) while Thought-in-the-Act supplies the operational logic (set the conditions, embody emergence, then iterate on technique). Practically speaking, we structure research as sequences of emergent experiments with practitioners: define initial conditions from Black sonic epistemologies, build minimal instruments that condition those relations, and let the ensuing practices “launch concepts-in-the-making” that are captured in reflective documentation, code, and performance \parencite[p.~89]{Manning_Massumi_2014}. In Loveless’s terms, this is a pedagogy of the lab and of the classroom. In hearkening to Haraway, a "common livable world must be composed, bit by bit, or not at all," hopefully resulting in new worlds of listening and making \parencite[p.~102]{Loveless_2019}.

